% !TeX program = pdflatex
\documentclass[12pt,a4paper]{report}

% Limbă și diacritice
\u23sepackage[romanian]{babel}
\usepackage[T1]{fontenc}
\usepackage[utf8]{inputenc}
\usepackage{csquotes}

% Margini & font
\usepackage[margin=2.5cm]{geometry}
\usepackage{lmodern}
\linespread{1.1}

% Hiperlinkuri
\usepackage[hidelinks]{hyperref}

% Tabele & figuri
\usepackage{graphicx}
\usepackage{float}
\usepackage{booktabs}
\usepackage{siunitx}
\sisetup{output-decimal-marker={,}}
\usepackage{subcaption}

% Matematică
\usepackage{amsmath, amssymb, amsthm}

% Enumerări
\usepackage{enumitem}

% Cod (opțiunea 1 - listings, ușor de compilat)
\usepackage{listings}
\lstset{
  basicstyle=\ttfamily\small,
  frame=single,
  numbers=left,
  numberstyle=\tiny,
  breaklines=true,
  tabsize=2,
  showstringspaces=false,
  language=Python
}

% (Opțiunea 2 - minted, comentată implicit; necesită -shell-escape)
% \usepackage{minted}
% \setminted{fontsize=\small, breaklines=true, frame=single}

% Numerotări: Section -> I., Subsection -> 1.1 etc.
\setcounter{secnumdepth}{3}
\setcounter{tocdepth}{3}

% Informații titlu (modifică după nevoie)
\newcommand{\Universitate}{UNIVERSITATEA DE STAT DIN MOLDOVA}
\newcommand{\Facultate}{FACULTATEA DE ECONOMIE}
\newcommand{\Departament}{DEPARTAMENTUL CONTABILITATE ȘI INFORMATICĂ ECONOMICĂ}
\newcommand{\TitluProiect}{Preferințe de cafea și consum zilnic}
\newcommand{\Disciplina}{Statistica pentru Știința Datelor}
\newcommand{\Autor}{Ivan Majeru}
\newcommand{\Titular}{Conf. univ. dr. Ludmila Novac}
\newcommand{\An}{2025}

\begin{document}

% Pagina de titlu
\begin{titlepage}
  \centering
  {\Large \Universitate\par}
  {\large \Facultate\par}
  {\large \Departament\par}
  \vspace{2cm}
  {\LARGE Cercetare statistică\par}
  \vspace{0.5cm}
  {\Large \TitluProiect\par}
  \vspace{2.5cm}
  \begin{flushleft}
    \textbf{Autor:} \Autor\par
    \textbf{Titular disciplină:} \Titular\par
    \textbf{Disciplina:} \Disciplina\par
  \end{flushleft}
  \vfill
  {\An\par}
\end{titlepage}

\pagenumbering{roman}
\tableofcontents
\listoffigures
\listoftables
\clearpage
\pagenumbering{arabic}

% I. OBSERVAȚIA STATISTICĂ
\section*{I. Observația statistică}
\addcontentsline{toc}{section}{I. Observația statistică}
\setcounter{section}{1}

\subsection{Scopul observației statistice}
Scopul cercetării este de a identifica preferințele de cafea (tipuri preferate, loc de achiziție) și de a analiza consumul zilnic (număr de căni) în funcție de caracteristici socio-demografice.

\subsection{Populația statistică}
Populația vizată include persoane adulte (18+) din Republica Moldova, rezidente în mediul urban și rural.

\subsection{Metoda de observație}
Anchetă statistică prin chestionar online/offline. Eșantionare neprobabilistică (de confort) sau probabilistică (dacă este cazul).

\subsection{Eșantionul observat}
Eșantion de 100 de persoane (date sintetice). Vârsta medie: 46{,}2 ani; consum mediu zilnic: 2{,}75 căni; buget mediu lunar: 1209{,}60 MDL.

\subsection{Caracteristici statistice}
Exemplu:
\begin{itemize}[leftmargin=1.2cm]
  \item Vârstă (ani)
  \item Sex (Masculin/Feminin)
  \item Tip localitate (Urban/Rural)
  \item Tip cafea preferat (Espresso, Latte, Cappuccino, Americano, etc.)
  \item Număr căni/zi
  \item Ora consum dominantă (Dimineața/Prânz/Seara)
  \item Buget lunar pentru cafea (MDL)
  \item Loc achiziție (Acasă-Boabe, Cafenea, Oficiu, etc.)
\end{itemize}

\subsection{Scale de măsurare}
\begin{itemize}[leftmargin=1.2cm]
  \item Vârstă, Număr căni/zi, Buget lunar: scară raport.
  \item Sex, Tip localitate, Tip cafea, Loc achiziție: nominal.
  \item Ora consum dominantă: ordinal (dacă o ordonezi), altfel nominal.
\end{itemize}

\subsection{Sursa datelor}
Date primare sintetice generate pentru scop didactic (Jupyter/Python).

% II. SISTEMATIZAREA ȘI PREZENTAREA DATELOR
\section*{II. Sistematizarea și prezentarea datelor observate}
\addcontentsline{toc}{section}{II. Sistematizarea și prezentarea datelor observate}
\setcounter{section}{2}

\subsection{Tabel descriptiv}
Introduceți tabelul cu primele observații și descrierea variabilelor.

\begin{table}[H]
\centering
\caption{Primele 5 observații din setul de date}
\label{tab:primele5}
\begin{tabular}{@{}lllllll@{}}
\toprule
ID & Vârstă & Sex & Localitate & Cafea & Căni/zi & Buget (MDL) \\
\midrule
1 & 67 & M & Urban & Espresso & 2 & 273 \\
2 & 22 & F & Urban & Latte & 3 & 1891 \\
3 & 47 & F & Urban & Cappuccino & 4 & 838 \\
4 & 41 & M & Urban & Cappuccino & 2 & 328 \\
5 & 20 & F & Rural & Americano & 4 & 1025 \\
\bottomrule
\end{tabular}
\end{table}

\subsection{Clasificarea unităților observate și ponderea pe grupe}
Exemple de tabele frecvențe/ponderi:
\begin{itemize}[leftmargin=1.2cm]
  \item După frecvența consumului (cani/zi grupate)
  \item După tipul de cafea preferat
  \item După locul de achiziție
  \item După sex / localitate
\end{itemize}

\begin{table}[H]
\centering
\caption{Structura după tipul de cafea preferat}
\label{tab:tip_cafea}
\begin{tabular}{@{}llllll@{}}
\toprule
Tip & Espresso & Latte & Cappuccino & Americano & Alte \\
\midrule
Frecvență &  &  &  &  &  \\
Pondere (\%) &  &  &  &  &  \\
\bottomrule
\end{tabular}
\end{table}

\subsection{Gruparea după caracteristici cantitative (intervale discrete/continue)}
Definiți intervalele folosind reguli de tip Sturges: 
$$
k = 1 + 3.322 \cdot \log_{10}(n), \quad h = \frac{x_{\max} - x_{\min}}{k}.
$$

\begin{table}[H]
\centering
\caption{Gruparea după numărul de căni pe zi (intervale continue)}
\label{tab:cani_intervale}
\begin{tabular}{@{}lllll@{}}
\toprule
Interval & [0,1) & [1,2) & [2,3) & [3,4) \\
\midrule
Frecvență &  &  &  &  \\
Pondere (\%) &  &  &  &  \\
\bottomrule
\end{tabular}
\end{table}

\subsection{Prezentare grafică a distribuțiilor și a structurii}
Includeți diagrame (histograme, bar chart, pie chart) exportate din Jupyter.

\begin{figure}[H]
  \centering
  \includegraphics[width=0.7\textwidth]{figuri/hist_cani_pe_zi.png}
  \caption{Distribuția numărului de căni pe zi}
  \label{fig:hist_cani}
\end{figure}

\begin{figure}[H]
  \centering
  \includegraphics[width=0.65\textwidth]{figuri/bar_tip_cafea.png}
  \caption{Structura după tipul de cafea preferat}
  \label{fig:bar_tip_cafea}
\end{figure}

% III. ANALIZA STATISTICĂ INFERENȚIALĂ
\section*{III. Analiza statistică inferențială folosind indicatori de tendință centrală}
\addcontentsline{toc}{section}{III. Analiza statistică inferențială folosind indicatori de tendință centrală}
\setcounter{section}{3}

\subsection{Indicatori de medie (pentru fiecare distribuție/grupare)}
Exemple (pentru Căni/zi):
$$
\bar{x} = \frac{1}{n}\sum_{i=1}^{n} x_i, \quad 
\tilde{x}=\text{mediană}, \quad
\text{media pătratică } = \sqrt{\frac{1}{n}\sum x_i^2}.
$$

\subsection{Indicatori de poziție (quartile, mediană, modă)}
\[
Q_1, \ Q_2(=\tilde{x}), \ Q_3 \quad \text{(pe baza distribuției ordonate sau a intervalelor)}